\chapter{Linear Solvers and PETSc}

For the committed 3D benchmark, the key linear-solver fact is that the outer
solver family and the inner preconditioner backend are now separate controls.
The current mainline combination is:
\begin{itemize}
\item \code{solver\_type = PETSC\_MATLAB\_DFGMRES\_HYPRE\_NULLSPACE},
\item \code{pc\_backend = pmg\_shell}.
\end{itemize}

So the historical solver-type name is preserved, but the actual backend used by
the default benchmark is the PMG shell path.

\section{Mainline Role Map}

\begin{longtable}{@{}p{0.29\textwidth}p{0.28\textwidth}p{0.35\textwidth}@{}}
\toprule
MATLAB role & MATLAB location & Mainline PETSc location \\
\midrule
\endhead
Solver factory & \matlabpkg{+LINEAR\_SOLVERS/set\_linear\_solver.m} & \petscmod{linear/solver.py::SolverFactory.create} \\
Deflated FGMRES logic & \matlabpkg{+LINEAR\_SOLVERS/DFGMRES.m} & MATLAB-style PETSc wrappers in \petscmod{linear/solver.py} \\
Iteration logging & \matlabpkg{+LINEAR\_SOLVERS/IterationCollector.m} & \petscmod{linear/collector.py} \\
Near-nullspace helpers & \matlabpkg{+LINEAR\_SOLVERS/near\_null\_space\_elasticity\_3D.m} & \petscmod{linear/preconditioners.py} and solver setup \\
\bottomrule
\end{longtable}

\section{What The Mainline Solver Actually Means}

For the default parallel PMG benchmark, the solver stack should be read as:
\begin{itemize}
\item \textbf{outer solver family:} PETSc-wrapped MATLAB-style deflated FGMRES,
\item \textbf{fine-grid operator:} the reduced free-space regularized operator
coming from the owned-row tangent path,
\item \textbf{preconditioner backend:} \code{pmg\_shell}.
\end{itemize}

That is the decisive interpretation. The benchmark is not currently using the
Hypre backend despite the historical text inside \code{solver\_type}.

\section{PMG Shell Requirements}

The PMG-shell path enforces a narrower contract than the broader solver
framework:
\begin{itemize}
\item \code{preconditioner\_matrix\_source = tangent},
\item \code{preconditioner\_matrix\_policy = current},
\item \code{preconditioner\_rebuild\_policy = every\_newton},
\item \code{full\_system\_preconditioner = false}.
\end{itemize}

Those requirements matter because they determine which constitutive and Newton
interfaces are actually exercised. In particular, they keep the mainline path on
the reduced free-space operator rather than the full-system or BDDC branches.

\section{Mainline Solver Lifecycle}

At each nonlinear step, the solver side does the following:
\begin{enumerate}
\item receive the current reduced operator from Newton,
\item decide whether the PMG preconditioner should be rebuilt,
\item configure the PMG shell multilevel objects,
\item solve the reduced linear system with the PETSc/MATLAB deflated FGMRES
wrapper,
\item record iteration counts, solve time, preconditioner time, and
orthogonalization time.
\end{enumerate}

This is the same broad job as the MATLAB solver object, but PETSc makes the
preconditioner lifecycle explicit.

\section{How The Mainline Path Relates To PMG}

The custom \code{pmg\_shell} backend is the current default PMG route for the
3D benchmark. In practical terms it means:
\begin{itemize}
\item the fine operator comes from the distributed owned-row tangent assembly,
\item the operator passed to the solve is the reduced free-space operator,
\item the preconditioner is rebuilt from the current tangent each Newton step.
\end{itemize}

That is the mainline meaning of ``parallel PMG'' in this repository today.

\section{Where To Edit}

\begin{itemize}
\item \petscmod{linear/solver.py} for solver-family selection, PMG-shell setup,
and preconditioner lifecycle,
\item \petscmod{linear/preconditioners.py} for helper objects and nullspace
construction,
\item \petscmod{linear/collector.py} for iteration and timing collection.
\end{itemize}

\section{Where The Other Solver Families Went}

Hypre, GAMG, BDDC, direct solves, built-in PETSc \code{pmg}, and other solver
policies are moved to Appendix~\ref{app:linear}.
